\documentclass{report}
\usepackage{amsmath,amssymb}
\setlength{\parindent}{0mm}
\setlength{\parskip}{1em}
\begin{document}
\begin{center}
\rule{6in}{1pt} \
{\large Stefan M. Wild \\
{\bf Optimal Finite Difference Derivatives of Noisy, Iterative Simulations}}

Mathematics and Computer Science Division \\ Argonne National Laboratory \\ 9700 S Cass Ave Bldg 240-1154 \\ Argonne IL 60439
\\
{\tt wild@mcs.anl.gov}\\
Jorge Mor\'e\end{center}

We have developed new mathematical theory and tools for estimating the
computational noise that arises in numerical simulations throughout
scientific computing
[4].
Roundoff errors, discretizations, numerical solutions to systems of
equations, and adaptive techniques are among the many sources
of computational noise, which can destroy the smoothness of
the processes underlying a simulation and complicate optimization,
sensitivity analysis, and other applications depending on the
simulation output.
Our theory is based on a stochastic model but does not
assume a specific form (e.g., Gaussian or mean zero)
for the distribution of the noise. Our technique is similar to that of
Hamming
[3]
and has been validated empirically on
deterministic simulations where the theory does not hold.

In this talk, we use an estimate of the computational noise to address a
longstanding problem in derivative estimation:
\emph{How should finite difference parameters be determined
when working with a noisy function?}

We have derived optimal finite difference parameters that are easy to
compute, depending only on an estimate of the computational noise and a
coarse bound on a higher-order derivative (typically requiring just a few
additional simulations)
[5].
Our estimates, $h_{\rm opt}$,
come with provable
approximation bounds for the resulting mean-squared error between the
finite difference estimate and the directional derivative of a smooth
function,
\begin{equation*}
\mathcal{E}(h_{\rm opt}) \leq \gamma \min \limits_{h\leq h_U}
\mathcal{E}(h), \qquad
\mathcal{E}(h) = \mathbb{E}
\left\{\left(\frac{f(x_0+hp)-f(x_0)}{h}-f'_s(x_0;p) \right)^2 \right\}.
\label{eq:bound}
\end{equation*}
An exciting aspect of this work is that we can obtain bounds on the number
of correct digits in the noisy derivative estimate. For example, for
forward differences, a typical approach is to use a multiple of the square
root of the machine's precision. Our numerical experiments show that in
many applications our estimate obtains 2-3 more digits of accuracy in the
derivative than does the classical approach.

Here we use the symmetric positive definite matrices
in the Florida Sparse Matrix collection
[2]
to
define noisy, deterministic quadratics of the form
\begin{equation*}
f(x) = \|y_{\tau}\|_2^2, \mbox{ where } Ay_{\tau}=x,
\end{equation*}
where $y_{\tau}$ is obtained using an iterative solver with tolerance
$\tau$.
We consider Krylov solvers in \texttt{MATLAB}
[1],
but the
numerical results are similar for other solvers.

We use these finite difference estimates to show how computational
noise can destroy the accuracy of derived calculations, for example, the
computation of derivatives.
In some cases, the accuracy of derivative estimates based on function
values is many times better than that of finite-precision evaluation of
derivatives obtained by hand-coded or automatically differentiated codes.
Our experiments also illustrate that the relationship between
truncation errors and the computational noise can be surprisingly
nonintuitive.

\paragraph{References}
\begin{enumerate}
\item[{[1]}] R. Barrett, M. Berry, T. F. Chan, J. Demmel, J.
Donato, J. Dongarra, V. Eijkhout, R. Pozo, C. Romine, and H. Van der Vorst.
{\em Templates for the Solution of Linear Systems: Building Blocks for
Iterative Methods, 2nd Edition}. SIAM, Philadelphia, PA, 1994.

\item[{[2]}] T. A. Davis and Y. F. Hu.
``The University of Florida Sparse Matrix Collection.''
\textit{ACM Transactions on Mathematical Software}, Vol.~38(1), pp.
1--25, (2011).
DOI \url{http://dx.doi.org/10.1145/2049662.2049663}

\item[{[3]}] R. W. Hamming. {\em Introduction to Applied
Numerical Analysis.} McGraw-Hill, 1971.

\item[{[4]}]
J.~Mor\'e and S. M.~Wild.
``Estimating Computational Noise.''
{\em SIAM J.~Scientific Computing}, Vol.~33(3), pp. 1292--1314, (2011). \\
DOI \url{http://dx.doi.org/10.1137/100786125}

\item[{[5]}]
J.~Mor\'e and S. M.~Wild.
``Estimating Derivatives of Noisy Simulations.''
\textit{ACM Transactions on Mathematical Software}, Vol.~38(3), to
appear, (2012).
URL \url{http://www.mcs.anl.gov/uploads/cels/papers/P1785.pdf}
\end{enumerate}


\end{document}
